
\documentclass[../../../physics-standard_answers_all.tex]{subfiles}

\begin{document}

\setlength{\abovedisplayskip}{2mm}
\setlength{\belowdisplayskip}{1mm}

\begin{enumerate}
    \item AB間を流れる電流の大きさを$I$とする．キルヒホッフ則より，
    %
    \begin{align*}
        3RI + 2RI = V_0 \qquad \text{∴} \; I = \frac{V_0}{5R} \;.
    \end{align*}
    %
    よって，
    \begin{align*}
        V_{2} = 3RI = \uwave{\frac{3}{5}V_0} \;.
    \end{align*}

    \setlength{\abovedisplayskip}{2mm}
    \setlength{\belowdisplayskip}{2mm}

    \item R$_1$の電流を$I$（右向き），R$_2$の電流を$i$（右向き）とすると，
    R$_3$の電流は$I - i$（右向き）である．キルヒホッフ則より，
    %
    \begin{align*}
        \left\{
        \begin{aligned}
            & Ri = 2R(I - i) \;, \\
            & Ri + 2RI = V_0
        \end{aligned}
        \right.
        \qquad \text{∴} \; i=\uwave{\frac{V_0}{4R}} \;,
        \quad I = \frac{3V_0}{8R} \;.
    \end{align*}

    \item R$_2$の電流を$i$（右向き）とすると，R$_3$の電流は$I-i$（右向き）である．
    キルヒホッフ則より，
    %
    \begin{align*}
        \left\{
        \begin{aligned}
            & Ri + RI = 2V_0 \;, \\
            & Ri = 2R(I - i) + V_0
        \end{aligned}
        \right.
        \qquad \text{∴} \; i=\frac{V_0}{R} \;,
        \quad I = \uwave{\frac{V_0}{R}} \;.
    \end{align*}

    \item R$_1$の電流を$i$（右向き），R$_5$の電流を$j$（上向き）とすると，
    R$_2$の電流は$i+j$（右向き），R$_3$の電流は$I-i$（右向き），
    R$_4$の電流は$I-i-j$（右向き）である．キルヒホッフ則より，
    %
    \begin{align*}
        \left\{
        \begin{aligned}
            & R(i+j) + 2Ri = R(I-i-j) + R(I-i) = V_0 \;, \\
            & 2Ri = Rj + R(I-i) \;.
        \end{aligned}
        \right.
    \end{align*}
    よって，
    \begin{align*}
        i = \frac{4V_0}{13R} \;, \quad j = \frac{V_0}{13R} \;,
        \quad I = \uwave{\frac{11V_0}{13R}} \;.
    \end{align*}

\end{enumerate}

\end{document}
